\chapter{Escenario 1: caso empresarial MediSalud}
\sourcebadge{Actividad desarrollada literalmente desde el Documento Padre, escenario 1.}

MediSalud Ecuador opera un HIS con procesos de agendamiento, admisión, historia clínica electrónica (HCE), prescripción, farmacia, facturación, telemedicina y reportes gerenciales. La problemática no es únicamente disponibilidad técnica: los usuarios encuentran lentitud, abandonos, duplicidades y riesgos que afectan objetivos reales.

\begin{table}[H]\centering
\caption{Matriz de análisis inicial}
\begin{tabularx}{\textwidth}{>{\bfseries}p{0.31\textwidth}X}
\toprule
Pregunta & Respuesta del equipo\\\midrule
Procesos más críticos & Atención y registro de HCE; facturación y seguros; agendamiento y admisión.\\
Usuarios más afectados & Médicos tratantes, pacientes y personal de admisión.\\
Evidencia disponible & Quejas, incidentes y disponibilidad de servidores.\\
Evidencia faltante & Éxito de tareas, tiempos reales, errores, satisfacción y variación por contexto.\\\bottomrule
\end{tabularx}
\end{table}

La evaluación se concentra en la experiencia observable y no utiliza métricas internas del código como sustituto de la calidad en uso.

